\documentclass{article}

\usepackage{standards}
\usepackage{tabularx}
\usepackage{algorithm}
\usepackage{algpseudocode}
\usepackage{float}      % for [H]
\usepackage{amsmath,amssymb}
\usepackage{longtable}
\usepackage{wrapfig}
\usepackage{multicol}


% Hide the printed "EndFor" keyword (but the \EndFor token is still required)
\algblockdefx[For]{For}{EndFor}[1]{\textbf{for} #1 \textbf{do}}{}


% Comments style (avoids math in comments)
\algrenewcommand\algorithmiccomment[1]{\hfill\texttt{// }#1}

\title{Transformers}
\author{Sam}
\date{}

\begin{document}
\maketitle
\tableofcontents

Puedo comenzar estadísticas? imporatante

tener para la siguiente meeting
y una tabla pequeña con las características de los modelos usados

Más de un benchmark por dataset? genial
Fijar tres tareas: clasificación, detección (bounding boxes), captioning, visual question answering

Está bien preguntar al modelo categorizar píxeles? puede devolverme en formato imagen?
NO HACER ESTO

ENFOCARME EN ESTO!!!

``Dataset scripts are no longer supported, but found VQAv2.py''

Cuantos models: 3 diferentes
seleccionar 3 y pq
comparar 5 (pueden ser menos)
rigurosidad  científica
tiempo de ejecucion en x, accuracy en y

comentar modalidades (esto debe de aparecer en hugging face)

martes 14 9:30am









\section{Datasets}

\subsection{Dataset selection procedure}

For each period $\{$pre-2015, 2015, 2016, \dots, 2024, 2025$\}$, I selected the three multimodal datasets with the highest number of citations on \texttt{Semantic Scholar} (\texttt{semanticscholar.org/}) as Google scholar does not allow me to order items by number of citations. Concretely, for each year bin I queried \texttt{Semantic Scholar} with the keywords \texttt{``dataset''}, \texttt{``image''}, and \texttt{``objects''} and then ranked results by the number of citations. A dataset was included only if it satisfied all of the following:

\begin{enumerate}
  \item \textbf{Search visibility.} The dataset must appear in \texttt{Semantic Scholar} search results for at least one of the keywords \emph{dataset}, \emph{image}, or \emph{objects}.
  \item \textbf{Multimodality with vision.} The dataset must be multimodal and include \textbf{images or videos}. Purely non-visual multimodal datasets are excluded (e.g., ESC (2015) is excluded because it focuses on audio). Datasets with 3D/geometry are acceptable \textbf{if} they also include 2D images or videos (e.g., FlyingThings3D (2015) qualifies); datasets lacking a 2D visual modality are excluded (e.g., Google Scanned Objects (2022) excluded).
  \item \textbf{Scope exclusions.} The dataset must \textbf{not} be:
    \begin{itemize}
      \item fully medical (e.g., CheXpert (2019) excluded),
      \item exclusively aerial/remote-sensing imagery (e.g., EuroSAT (2017) excluded),
      \item images or videos taken from vehicles' cameras (e.g., nuScenes (2019) excluded;\, HDR cameras on top of vehicles such as the ones used for the street view on Google Maps are OK, hence Cityscapes (2016) is included).
    \end{itemize}
\end{enumerate}

Applying these rules, I retained, for each period, the top three qualifying datasets by publication count.




































































\subsection{Summary of the Datasets as a Whole}


\begin{longtable}{p{2.7cm}p{1.7cm}p{4cm}p{4cm}}
\caption{Summary of major multimodal / vision datasets (2009–2025). Terminology is explained in detail in the \textbf{Summary of Datasets as a Whole} section.}
\label{tab:datasets} \\
\hline
\textbf{Name} & \textbf{Citations} & \textbf{License} & \textbf{Type of elements}\\
\hline
\endfirsthead
\multicolumn{4}{l}{\textit{(continued from previous page)}}\\
\hline
\textbf{Name} & \textbf{Citations (approx.)} & \textbf{License} & \textbf{Type of elements}\\
\hline
\endhead
\hline \multicolumn{4}{r}{\textit{Continued on next page}}\\
\endfoot
\hline
\endlastfoot

ImageNet (2009) & 66,633 & Free for non-commercial research & Images + labels\\
\hline
UCF101 (2012) & 6,460 & CC0 1.0 Universal & Videos + action labels\\
\hline
MS COCO (2014) & 46,361 & CC BY 4.0 & Images + object segmentation + captions\\
\hline

Flickr30k Entities (2015) & 2,226 & Free for non-commercial research & Images + captions + bounding boxes + relationships\\
\hline
FlyingThings3D (2015) & 2,758 & Free for non-commercial research & Images (pairs) + 3D data\\
\hline
LSUN (2015) & 2,436 & Free for non-commercial research + "certain condicions" & Images + labels\\
\hline

Cityscapes (2016) & 12,189 & Free for non-commercial research & Images + segmentation + labels\\
\hline
Visual Genome (2016) & 6,018 & CC BY 4.0* & Images + bounding boxes + relationships + QA pairs\\
\hline
VQA v2.0 (2016) & 3,558 & CC BY 4.0 & Images (paired) + QA pairs\\
\hline

Fashion-MNIST (2017) & 9,435 & MIT License & Images + labels\\
\hline
Kinetics Human Action Video Dataset (2017) & 3,981 & CC BY 4.0* & Videos + labels\\
\hline
iNaturalist (2017) & 1,763 & CC BY-NC (Creative Commons Non-Commercial) & Images + boudning boxes + labels\\
\hline

Places (2018) & 4,319 & Free for non-commercial research purposes & Images + labels\\
\hline
Conceptual Captions (2018) & 2,643 & Free for any purpose & Images + captions\\
\hline
Open Images V4 (2018) & 1,441 & CC BY 2.0 (images) / CC BY 4.0 (annotations) & Images + bounding boxes + labels + relationships\\
\hline

GQA (2019) & 2,402 & CC BY 4.0 & Images + QA pairs\\
\hline
MVTec AD (2019) & 1,549 & CC BY-NC-SA 4.0 & Images + segmentation\\
\hline
LVIS (2019) & 1,492 & CC BY 4.0 & Images + segmentation + labels\\
\hline

DocVQA (2020) & 930 & ? & Images + QA pairs\\
\hline
DFDC (2020) & 597 & Meta research license? & Videos + labels\\
\hline
TextCaps (2020) & 469 & CC BY 4.0 & Images + captions\\
\hline

LAION-400M (2021) & 1,486 & CC BY 4.0 (metadata); images under original copyright & Image + labels + scores\\
\hline
FairFace (2021) & 616 & CC BY 4.0 & Images + labels\\
\hline
HDTF (2021) & 413 & GNU GPLv3 & Videos + audio + spoken text\\
\hline

LAION-5B (2022) & 4,059 & CC BY 4.0 (metadata); images under original copyright & Images + captions + scores\\
\hline
DiffusionDB (2022) & 349 & CC0 1.0 (data) / MIT (code) & Images + prompts + scores\\
\hline
TAD66K (2022) & 111 & Apache License 2.0 & Images + scores\\
\hline

MagicBrush (2023) & 391 & CC BY-SA 4.0 & Images + prompts\\
\hline
Pick-a-Pic (2023) & 559 & ? & Images + prompts + preference\\
\hline
InternVid (2023) & 411 & CC BY-NC-SA 4.0 & Videos + captions\\
\hline

OpenVid-1M (2024) & 149 & ? & Videos + prompts + captions\\
\hline
Visual CoT (2024) & 148 & CC BY-SA 4.0 & Images + bounding boxes + QA pairs\\
\hline
HQ-Edit (2024) & 104 & CC BY-NC-4.0 & Images + prompts\\
\hline

BLIP3o-60k (2025) & 91 & Apache License 2.0 & Images + prompts\\
\hline
ImgEdit (2025) & 38 & Apache License 2.0 & Images + prompts\\
\hline
ShareGPT-4o-Image (2025) & 27 & Apache License 2.0 & Images + prompts\\
\hline
\end{longtable}


In Table~\ref{tab:datasets}, when images are grouped or classified into categories, we count this as having \bfit{labels}. Labels are category tags that describe what the image or video shows. For example, a picture of a dog in \textbf{ImageNet} has labels such as \texttt{dog}, \texttt{mammal}, \texttt{animal}, while a picture of a woman laughing in \textbf{UCF101} is labeled as \texttt{smiling}.  \bfit{Captions} go a step further by describing what is happening or what can be seen in the image or video, as seen in Figure~\ref{fig:conceptual-captions} for an image captioned "\texttt{pop artist performs at the festival in a city}" in the \textbf{Conceptual Captions} dataset.\bfit{Prompts} are pieces of text used to generate or edit images in datasets built around generative models. Examples can be seen in Figures~\ref{fig:diffusiondb} or \ref{fig:magicbrush} for the \textbf{DiffusionDB} or \textbf{MagicBrush} datasets respectively.\\

\bfit{Bounding boxes} pinpoint specific object instances within an image by drawing rectangles around them (using coordinates such as $(x_{\min},y_{\min},x_{\max},y_{\max})$). See Figures \ref{fig:flickr30k} or \ref{fig:visual-genome} from \textbf{Flickr30K} or \textbf{Visual Genome} respectively as examples. \bfit{Segmentation} goes a step further by outlining the exact shape of each object, assigning every pixel to a class or instance. Semantic segmentation labels each pixel by class, while instance segmentation gives each object its own mask. See Figures \ref{fig:mscoco_ucf101} or \ref{fig:cityscapes} from \textbf{MS COCO} or \textbf{Cityscapes} respectively for examples.\\

Finally, \textbf{Pick-a-Pic} captures human \bfit{preferences} for images generated by text-to-image diffusion models. Images in the dataset come in pairs, with one of the two image being preferred. Figure~\ref{fig:pickapic2} shows examples where the chosen image is highlighted and the rejected one is darkened. The term \bfit{scores} refers to any numerical metadata. For example, an aesthetic score as in \textbf{TAD66K} or an NSFW score as in \textbf{LAION-5B}. \bfit{QA-pairs} (question–answer pairs) are exactly what they sound like: a question about the image and its corresponding answer, such as “\texttt{What type of fruit in the image is round?}” and “\texttt{apple}”, as shown in Figure \ref{fig:gqa} for the \textbf{GQA} dataset.\\


 \bfit{Relationships} refer to datasets that explicitly describe how objects interact or relate to one another, rather than just listing them. Only two datasets meet this criterion: \textbf{Flickr30K Entities} and \textbf{Visual Genome}. The former connects phrases in captions that refer to the same entity (like “\texttt{a boy}” and “\texttt{the child}”) and links them to specific bounding boxes in the image, as shown in Figure~\ref{fig:flickr30k}. The latter represents scenes as graphs of objects and their relationships e.g., “\texttt{man \textit{riding} horse}”, as shown in Figure~\ref{fig:visual-genome}. Finally, when a dataset contains \bfit{3D data}, it includes information about the geometry of the scene beyond standard RGB images. This only applies to \textbf{FlyingThings3D}, which contains stereo image pairs (left and right views rendered from two virtual cameras) along with detailed geometric information such as depth (disparity) maps, optical flow, segmentation masks, material indices, motion boundaries, and camera calibration data. Together, these describe the full 3D structure and motion of the scene, as shown in Figure~\ref{fig:flyingthings3d}.\\

Regarding the licenses, 





































\subsection{Summary of each Datasets}

\begin{multicols}{2}
\subsection{The Datasets}

\begin{figure*}[t]
    \centering
    \includegraphics[width=\textwidth]{ImageNet.jpg}
    \caption{Example of ImageNet hierarchical organization, showing a path from \textit{mammal} to \textit{husky}.}
\end{figure*}

\textbf{ImageNet} is a large-scale database of human-annotated photographs organized into a hierarchy of visual concepts. Each category, or “synset,” groups together images of the same concept, such as "husky", and is linked to related categories within a larger hierarchy, such as "animal".  
At the time of publication, ImageNet contained 12 major groups, covering 5,247 categories and about 3.2 million images.\\

















\begin{figure}[H]
    \centering
    \vspace{-10pt}
    \includegraphics[width=0.5\textwidth]{UCF101.jpg}
    \caption{Sample frames from UCF101 showing diverse action classes such as \textit{Apply Eye Makeup}, \textit{Playing Dhol}, and \textit{Baby Crawling}.}
    \vspace{-10pt}
    \label{fig:mscoco_ucf101}
\end{figure}

\textbf{UCF101 (2012).} The University of Central Florida (UCF) compiled a total of 13,320 video clips amounting to about 27 hours of footage. Each clip was then labeled as one of 101 action classes, 51 of which preserved audio. The videos are stored at a resolution of 320×240 pixels and 25 frames per second. On average, clips are about 7 seconds long.\\ 



















\begin{figure}[H]
    \centering
    \vspace{-10pt}
    \includegraphics[width=0.5\textwidth]{MSCOCO.jpg}
    \caption{Examples from the MS COCO dataset showing object segmentations for categories such as \textit{cow}, \textit{train}, and \textit{car}}
    \vspace{-10pt}
    \label{fig:mscoco_ucf101}
\end{figure}


\textbf{Microsoft COCO (2014).} The Microsoft COCO, abbreviated MS COCO, dataset contains photographs of everyday environments filled with multiple objects in their natural context. Each object instance is carefully annotated with detailed segmentations, allowing precise localization within the image as shown in Figure ~\ref{fig:mscoco_ucf101}.

MS COCO includes 91 common object categories that are easily recognizable, with over 2.5 million labeled instances across 328,000 images.\\











































\begin{figure}[H]
    \centering
    \vspace{-10pt}
    \includegraphics[width=0.35\textwidth]{Flickr30k1.jpg}
    \caption{Example image with captions; phrases in the same color are coreferent and linked to their corresponding bounding boxes.}
    \label{fig:flickr30k}
    \vspace{-10pt}
\end{figure}

\textbf{Flickr30k Entities} extends the Flickr30k dataset by providing detailed links between image regions and textual phrases. With 31{,}783 images and 158{,}915 captions, the dataset connects mentions that refer to the same entity across captions describing the same image, forming about 244{,}000 coreference chains, and associates these with roughly 276{,}000 manually drawn bounding boxes. An example of the final annotated data is shown in Figure~\ref{fig:flickr30k}, where color-coded phrases correspond to specific image regions.  These annotations enable the new task of \emph{phrase localization}: given an image and a caption, predict the region that corresponds to a particular entity mention.  
























\begin{figure}[H]
    \centering
    \vspace{-10pt}
    \includegraphics[width=0.5\textwidth]{FlyingThings3D2.jpg}
    \caption{Example scenes from the \textit{FlyingThings3D} dataset. The third row shows optical flow images, the fourth row shows disparity images, and the fifth row shows disparity change images (normalized for display).}
    \label{fig:flyingthings3d}
    \vspace{-10pt}
\end{figure}


\textbf{FlyingThings3D}~\cite{mayer2016large} consists of roughly 25,000 stereo frames rendered using \textit{Blender}, each displaying a  random arrangement of textured 3D objects. Each frame provides left and right RGB images, together with corresponding maps for optical flow, disparity, and disparity change, as well as object segmentation masks, material indices, and motion boundaries as shown in Figure~\ref{fig:flyingthings3d}. Complete camera calibration data and 3D geometry information are also provided.  \\














\textbf{LSUN} is a very large dataset consisting of about 69 million images belonging to 30 categories. Only a small portion of the images were labeled by humans, the rest being labeled by a trained classifier (e.g., for the \emph{train} category, less than $1/60$ of images were labeled by people), achieving ~90\% precision. Training popular convolutional networks on this larger, if noisier, data leads to noticeable gains; for example, fine-tuning with LSUN reduced a reported classification error from 28.6\% to 22.2\% on a related evaluation.\\





























\begin{figure*}[t]
    \centering
    \includegraphics[width=\textwidth]{Cityscapes.jpg}
    \caption{Example scenes from the \textit{Cityscapes} dataset. From left to right: fine annotations for training/validation (3,475 images), coarse annotations for training (20,000 images), and fine annotations for testing (1,525 images). Each image shows dense semantic and instance-level labeling of complex urban street scenes.}
    \label{fig:cityscapes}
\end{figure*}

\textbf{Citiscapes} contains stereo video sequences recorded across 50 different cities, primarily in Germany. From these recordings, 5,000 images received high-quality fine annotations, while an additional 20,000 were annotated coarsely (less boundary precision) with a ~98\% accuracy after excluding ambiguous (void) regions. Figure~\ref{fig:cityscapes} illustrates examples of both fine and coarse annotations used for training, validation, and testing. The dataset is split into 2,975 training, 500 validation, and 1,525 test images, ensuring balance across geography, population size, and season. In total, 30 semantic classes (19 for evaluation) are grouped into eight high-level categories: flat, construction, nature, vehicle, sky, object, human, and void.  
































\begin{figure*}[t]
    \centering
    \includegraphics[width=\textwidth]{VisualGenome1.jpg}
    \caption{Example image and scene graph from the \textbf{Visual Genome} dataset. The left side shows labeled objects with bounding boxes, while the right side illustrates the relationships between these objects in a graph.}
    \label{fig:visual-genome}
\end{figure*}

\textbf{Visual Genome}~\cite{krishna2017visualgenome} human-annotated imagea describe not only \emph{what} objects are present, but also \emph{how} they interact. The dataset contains over 100,000 real-world images, each annotated with an average of 21 objects, 18 attributes, and 18 pairwise relationships, as exemplified in Figure ~\ref{fig:visual-genome}. Beyond scene graphs, Visual Genome includes over 1.7 million question–answer pairs designed to test visual reasoning and comprehension. 






























\begin{figure}[H]
    \centering
    \vspace{-10pt}
    \includegraphics[width=0.5\textwidth]{VQAV2.jpg}
    \caption{Examples from the balanced \textbf{VQA v2.0} dataset. Each question is paired with images that are visually similar yet lead to different answers, encouraging models to use visual evidence rather than language priors.}
    \label{fig:vqa-v2}
    \vspace{-10pt}
\end{figure}


\noindent
\textbf{VQA v2.0.} For each image in the Visual Question Answering (VQA) dataset, annotators tried to find a complementary image where the question still makes sense, but has a different answer. These image pairs with question-answers result in VQA v2.0, as shown in Figure ~\ref{fig:vqa-v2}. In total, they collected approximately 195K complementary images for train, 93K complementary images for val, and 191K complementary images for test set.

Benchmarking shows that state-of-the-art models trained on the original data perform significantly worse on the balanced set; clear evidence that they had exploited language regularities. Training on the balanced data improves results but the task remains challenging. 






























\begin{figure*}[t]
    \centering
    \includegraphics[width=\textwidth]{FASHION-MNIST.png}
    \caption{The conversion pipeline used to generate the \textit{Fashion-MNIST} dataset. Each row shows an example from a product category (dress, sandal), with transformations applied sequentially until the final (rightmost) image that will be part of the dataset is created.}
    \label{fig:fashion-mnist}
\end{figure*}

\textbf{Fashion-MNIST}~\cite{xiao2017fashionmnist} consists of 70,000 grayscale images of size $28\times28$, evenly distributed across 10 fashion categories such as T-shirts, trousers, coats, and shoes. The training set contains 60,000 images, and the test set includes 10,000 images. As such, the dataset is similar in structure to the MNIST digits dataset, from where it gets its name. Fashion-MNIST poses a significantly harder classification challenge: typical models that achieve over 99.7\% accuracy on digits reach only around 89–90\% on Fashion-MNIST.\\



























\begin{figure*}[t]
    \centering
    \includegraphics[width=\textwidth]{KHAVD.png}
    \caption{Example frames from the ``brushing hair'' class in the \textbf{Kinetics Human Action Video Dataset}, illustrating variation in appearance, lighting, and camera viewpoint. Each clip corresponds to a unique YouTube video.}
    \label{fig:kinetics-brushing}
\end{figure*}

The \textbf{Kinetics Human Action Video Dataset} consists of 400 action classes with approximately 400–1150 ten-second clips per class, each extracted from a distinct YouTube video (for a total of $\sim$306K clips). Actions span three major categories: \emph{person-only actions} (e.g., laughing, brushing hair), \emph{person–object interactions} (e.g., playing instruments), and \emph{person–person interactions} (e.g., shaking hands). Clips retain audio but are designed for visual classification, enabling use in multimodal learning tasks. The dataset is divided into three subsets with 250–1000 training clips per class, 50 validation clips per class, and 100 test clips per class.\\






























\begin{figure}[H]
    \centering
    \vspace{-10pt}
    \includegraphics[width=0.5\textwidth]{iNaturalist.jpg}
    \caption{Sample bounding-box annotations for the \textbf{iNaturalist} dataset. Annotators drew boxes for the \emph{super-class} (e.g., “box all birds”), up to ten instances per image; boxes could overlap, cover only visible parts when occluded, and a single box could enclose physically connected instances.}
    \label{fig:inat2017-fig3}
    \vspace{-10pt}
\end{figure}

\textbf{iNaturalist (2017)}, at time of publication, contained 859,000 images (579,184 train, 95,986 val, and 182,707 test images) spanning 5,089 plant and animal species. Labels reflect the consensus of informed contributors on the iNaturalist platform, and annotators drew boxes around the plant/animals that appear on the image. On classes with the most training examples, the best classifiers at the time could only reach 67\% accuracy.\\


















\begin{figure*}[t]
    \centering
    \includegraphics[width=\textwidth]{Places.jpg}
    \caption{Example scenes from the \textit{Places} dataset, illustrating variations across categories such as bedrooms, kitchens, forests, and coasts. Each scene combines structural and stylistic diversity captured through adjective-based search expansions.}
    \label{fig:places}
\end{figure*}

The \textbf{Places} database~\cite{zhou2017places} comprises 10 million images labeled with 434 semantic scene categories (e.g., "messy kitchen", "misty coast"). Designed to complement object-centric datasets like ImageNet, it provides exhaustive coverage of environments such as \textit{beaches}, \textit{corridors}, \textit{kitchens}, and \textit{forests}, enabling models to understand not just objects, but the broader spatial and functional context in which they appear. An unusual and creative element of the collection process was the use of 696 English adjectives to diversify visual queries — combinations such as “messy kitchen,” “romantic bedroom,” and “misty coast.”


















\begin{figure}[H]
    \centering
    \vspace{-10pt}
    \includegraphics[width=0.5\textwidth]{ConceptualCaptions.jpg}
    \caption{Examples from the \textbf{Conceptual Captions} dataset. Original web \texttt{alt}-texts (top) are automatically cleaned and hypernymed into concise captions (bottom).}
    \label{fig:conceptual-captions}
    \vspace{-10pt}
\end{figure}

The \textbf{Conceptual Captions} dataset~\cite{sharma2018conceptual} contains about 3.3 million image–caption pairs automatically harvested from web \texttt{alt}-texts, making it over ten times larger than \textbf{MS COCO}. The final split includes 3.3 million training examples, 28,000 validation, and 22,500 test pairs, with captions averaging ten words in length. 

The dataset was created through a multi-stage automatic pipeline that processed billions of webpages. The first stage filtered out non-JPEG, small, distorted, or offensive images. In the second stage, noisy or incomplete \texttt{alt}-texts—such as keyword spam or hashtag lists—were discarded using part-of-speech, sentiment, and profanity analysis. Next, candidate captions were retained only if there was lexical overlap between concepts detected in the image (via Google Vision labels) and in the text. Finally, named entities, locations, and dates were replaced with hypernyms (e.g., “Harrison Ford” replaced by “actor”) to improve generalization. Human evaluation showed that over 90\% of captions were judged as good. 























\begin{figure*}[t]
    \centering
    \includegraphics[width=\textwidth]{OpenImagesV4.jpg}
    \caption{Example annotations from the \textit{Open Images V4} dataset for image classification, object detection, and visual relationship detection. Green labels mark present classes, red labels absent ones. Dashed boxes group pairs of objects sharing a visual relationship.}
    \label{fig:openimages}
\end{figure*}

The \textbf{Open Images V4} dataset~\cite{kuznetsova2020openimages} contains 9.2 million Creative Commons–licensed Flickr images with unified annotations for image classification, object detection, and visual relationship detection (see Figure~\ref{fig:openimages}). It is split into 9.0M training, 41K validation, and 125K test images. It provides 30.1 million image-level labels for 19.8k concepts, 15.4 million bounding boxes for 600 object classes, and 375k relationship annotations involving 57 classes. Each image contains on average eight annotated objects, reflecting the complexity of real-world scenes.

Compared to previous datasets such as MS-COCO or ImageNet, Open Images V4 offers more than 15 times as many bounding boxes and integrates multiple annotation types within the same images.\\




































\begin{figure}[H]
    \centering
    \includegraphics[width=.5\textwidth]{GQA.jpg}
    \caption{Example GQA images and queries illustrating spatial reasoning, attributes, and compositional question answering.}
    \label{fig:gqa}
\end{figure}

The \textbf{GQA} dataset~\cite{hudson2019gqa} is built upon the \textbf{Visual Genome} scene graphs and contains 113k images with 22 million automatically generated yet semantically grounded questions, of which 1.7 million form a balanced training and evaluation split. Each question is associated with a \textit{functional program} (a symbolic representation of its reasoning steps such as \texttt{select$\rightarrow$filter$\rightarrow$relate$\rightarrow$query}) and with visual justifications linking words to specific image regions.

A distinctive feature of GQA is its emphasis on compositionality: a single question may reference nested visual relationships such as “the apple on the table to the left of the bowl.” Human performance reaches 89.3\%, while strong VQA systems achieve only around 54.1\%, and a blind LSTM just 42.1\%, illustrating the task’s difficulty and the gap between recognition and reasoning. As the authors aptly remark, “It takes more than a smart guess to answer a good question,” encapsulating the spirit of the GQA benchmark.

























\begin{figure}[H]
    \centering
    \includegraphics[width=.5\textwidth]{MVTec AD.jpg}
    \caption{Examples from \textbf{MVTec AD}: each triplet shows a defect-free image (left), an anomalous image (middle), and the pixel-precise ground-truth mask overlaid on a zoomed region (right).}
    \label{fig:mvtec-ad}
\end{figure}

The \textbf{MVTec Anomaly Detection (MVTec AD)} dataset~\cite{bergmann2019mvtec}, created by the MVTec company, attempst to solve the issue that other datasets contain only normal (defect-free) imagery is available for training. It contains 5,354 high-resolution RGB/grayscale images spanning 15 categories (10 objects and 5 textures). Training splits comprise only normal images (3,629 total), while the test split (1,725 images) mixes normal and defective cases. Anomalies cover more than 73 defect types (e.g., scratches, dents, contaminations, missing parts, and structural deformations) and every defective image is annotated with pixel-accurate ground-truth regions (1,888 masks). Images were captured with industrial telecentric optics at $\approx700\!\times\!700$–$1024\!\times\!1024$ px (this means magnification of the image does not change with the object’s distance from the lens), though some categories are grayscale only.


































\begin{figure}[H]
    \centering
    \includegraphics[width=.5\textwidth]{LVIS2.jpg}
    \caption{Example annotations from the \textbf{LVIS} dataset, illustrating fine-grained segmentation masks across diverse kitchen scenes.}
    \label{fig:lvis}
\end{figure}

The \textbf{LVIS} dataset~\cite{gupta2019lvis} (\textit{Large Vocabulary Instance Segmentation}) introduces over 1,000 entry-level object categories and aims to collect approximately two million high-quality segmentation masks over 164,000 images. Built with an “evaluation-first” philosophy, LVIS retains compatibility with the COCO average precision (AP) metric but dramatically increases category diversity while embracing the natural long-tailed distribution of real-world objects.

Unlike previous benchmarks, which exhaustively annotate a fixed set of categories, LVIS adopts a \textit{federated} dataset design. Each category has its own positive and negative subsets of images that are exhaustively annotated for that specific class, avoiding the impractical burden of labeling every image with all categories. This federated approach ensures fairness during evaluation even when an object can belong to multiple overlapping or synonymous categories. As the authors explain, “a toy deer is both a toy and a deer,” and the benchmark must fairly credit detectors that output either or both labels. Notably, LVIS segments were found to have higher boundary accuracy and annotator consistency than COCO and ADE20K masks, achieving over 99\% acceptance after verification.






































\begin{figure}[H]
  \centering
  \includegraphics[width=0.5\textwidth]{DocVQA.jpg}
  \caption{\textbf{DocVQA example.} Sample document image with Q\&A pairs that require reading text and interpreting page layout/structure.}
  \label{fig:docvqa}
\end{figure}


\textbf{DocVQA} is a large-scale dataset for \emph{visual question answering on document images}. It targets purpose-driven understanding of real documents (letters, forms, tables, lists, figures, and handwritten notes) to train systems that must both \emph{read} text and \emph{reason} over layout and structure. The dataset contains \textbf{12{,}767} pages drawn from \textbf{6{,}071} industry documents (spanning early 1900s–2018 and multiple domains) and \textbf{50{,}000} natural-language questions with one or more extractive answers.


Questions are answerable by copying spans \emph{verbatim} from the page (an extractive QA setting). Success therefore couples robust OCR with structural understanding: models must leverage cues like headings, seals, key–value pairs, table cells, and relative position (e.g. “at the top left”, “in the seal”, “ZIP code”). Data are split \textbf{80/10/10} into train/val/test: 39,463/5,349/5,188 questions over 10,194/1,286/1,287 images. Human accuracy is reported at \textbf{94.36\%}, leaving a substantial gap for current VQA/MRC baselines that struggle especially on structure-heavy queries.





















\begin{figure*}[t]
    \centering
    \includegraphics[width=\textwidth]{DFDC.jpg}
    \caption{\textbf{DFDC face mosaic.} Example frames from the DFDC dataset. All identities are consenting paid actors recorded specifically for creating and detecting face swaps.}
    \label{fig:dfdc}
\end{figure*}

The \textbf{DeepFake Detection Challenge (DFDC)} dataset is a large, consented face-swap corpus designed for training and benchmarking Deepfake detectors. It comprises 100k+ 10-second clips sourced from 3,426 paid actors (average 14.4 videos/subject; mostly 1080p), totaling 38.4 days of footage and 25TB of raw data. Relative to previous datasets, DFDC is orders-of-magnitude larger, higher quality, and ethically collected with explicit consent.

Multiple techniques (Deepfake Autoencoder, morphable-mask, Neural Talking Heads, and so on) were used to make the fake videos so that detection models wouldn’t just overfit to one kind of fake, but would learn to generalize to many.\\




































\begin{figure*}[t]
    \centering
    \includegraphics[width=\textwidth]{TextCaps3.jpg}
    \caption{Examples from TextCaps demonstrating how captions combine visual description with textual information in the image. Bold phrases indicate paraphrasing or reasoning; underlined ones are directly copied OCR text.}
    \label{fig:textcaps2}
\end{figure*}

The \textbf{TextCaps} dataset~\cite{sidorov2020textcaps} introduces the task of image captioning with reading comprehension, where models must both recognize scene text and reason about it within the visual context. For example, recognizing a “red sign” is insufficient, the model must read “Mornington Crescent” to identify the location, or understand that “Rice is winning” implies interpreting a scoreboard rather than simply detecting numbers. The dataset contains 145,000 captions over 28,000 images belonging originally to OpenImages v3.























\begin{figure*}[t]
    \centering
    \includegraphics[width=\textwidth]{laion-400m.jpg}
    \caption{Example image retrieval results for the query ``cat with blue eyes'' using the LAION-400M web demo. The dataset supports large-scale multimodal similarity search using CLIP embeddings.}
    \label{fig:laion400m}
\end{figure*}

The \textbf{LAION-400M} dataset contains an astounding 400 million image–text pairs, filtered using \textbf{CLIP} to ensure semantic alignment between images and captions. It was created as a community-driven effort to provide a large-scale, public alternative to proprietary datasets used to train models such as CLIP and DALL·E. Each entry includes an image URL, text caption, CLIP embeddings for both modalities, metadata (license type, image size, and NSFW tag), and precomputed kNN indices for efficient similarity search. 




















\begin{figure*}[t]
    \centering
    \includegraphics[width=\textwidth]{FairFace.jpg}
    \caption{Example faces from (a) FairFace and (b–d) popular prior datasets. FairFace is designed to be racially balanced, unlike many existing sets that are dominated by White faces.}
    \label{fig:fairface_grid}
\end{figure*}

\textbf{FairFace} introduces 108,501 in-the-wild face images with labels for race, gender, and age, constructed to be balanced over seven race groups: White, Black, Indian, East Asian, Southeast Asian, Middle Eastern, and Latino, which were annotated manually. 

The paper highlights that many widely used face datasets (e.g., celebrity or politician collections) are skewed toward lighter-skin/White faces, which leads to uneven accuracy across demographics. Models trained on FairFace yield higher and more consistent accuracy across race and gender on these novel sets, even when FairFace’s training subset is down-sampled. 
























\begin{figure}[H]
    \centering
    \includegraphics[width=0.5\textwidth]{hdtf.jpg}
    \caption{Overview of one-shot high-resolution talking face generation from the HDTF dataset. Given a reference image and driving audio, the model synthesizes high-resolution videos with synchronized facial animation.}
    \label{fig:hdtf}
\end{figure}

\noindent
\textbf{HDTF}'s authors note that prior systems top out around 256$\times$256, so that they introduce High-Definition Talking Face (HDTF), a high-definition, in-the-wild dataset ($\sim$16 hours of 720p--1080p video, 300+ subjects, 10k+ sentences) curated from recent YouTube content so the crops remain sharp after 512$\times$512 resizing. They synthesize high-definition, lip-synchronous videos (demonstrations include portraits and paintings) and report both objective and user-study gains over landmark-based baselines, especially at $512\times 512$.
































\begin{figure*}[H]
    \centering
    \includegraphics[width=\textwidth]{laion-5b.jpg}
    \caption{Nearest-neighbor examples from LAION-5B using CLIP embeddings. For each text query (Q), the corresponding top retrieved image and its caption (C) are shown.}
\end{figure*}

The LAION-5B dataset is the largest openly available image–text dataset ever released, designed to democratize research on large-scale multimodal models such as \textbf{CLIP}, \textbf{GLIDE}, and \textbf{Stable Diffusion}. It contains 5.85 billion CLIP-filtered image–text pairs collected from Common Crawl, including 2.32 billion English pairs, 2.26 billion multilingual pairs covering over a hundred languages, and 1.27 billion language-agnostic samples, which often describe products or places. Each entry is accompanied by metadata such as the image URL, caption text, image dimensions, cosine similarity between the image and text embeddings, and safety-related scores like watermark and NSFW probabilities.

The accompanying web interface allows interactive dataset exploration, nearest-neighbor search via CLIP embeddings, and flexible subset generation. The dataset has been validated through extensive experiments. Models trained on \textbf{LAION-400M} reproduce the zero-shot classification performance of OpenAI’s CLIP on ImageNet and other benchmarks, confirming that the filtering and scale of LAION-5B provide data of comparable quality. Subsets such as LAION-Aesthetics and LAION-High-Resolution have also been derived to train models like Stable Diffusion, which achieve state-of-the-art results in text-to-image generation. 
























\begin{figure}[H]
    \centering
    \includegraphics[width=0.35\textwidth]{diffusiondb2.jpg}
    \caption{Examples from \textbf{DiffusionDB}, showing generated images with their corresponding prompts}
    \label{fig:diffusiondb}
\end{figure}

\noindent
\textbf{DiffusionDB} is the first large-scale, community-generated prompt dataset for text-to-image diffusion models. It contains over 14 million images, 1.8 million unique text prompts, and detailed \textbf{Stable Diffusion hyperparameters} including seed, step count, CFG scale, sampler type, image size, and timestamp. The dataset was collected from the official Stable Diffusion public Discord server, where users shared their generation results. To ensure safe and ethical use, the authors applied state-of-the-art NSFW detection models on both images and prompts, producing continuous safety scores that allow flexible filtering. Prompts were also analyzed linguistically and semantically: most are short (6–12 tokens) and in English, but the dataset spans over 30 languages.\\



















































\begin{figure}[H]
  \centering
  \includegraphics[width=0.45\textwidth]{tad66k.jpg}
  \caption{Examples from the \textbf{TAD66K} dataset. Each pair shows images from the same theme with their average aesthetic scores (in red) and predicted scores in parentheses. Even when scores are similar, different themes correspond to distinct aesthetic criteria.}
\end{figure}

The \textbf{TAD66K} dataset (Theme and Aesthetics Dataset) was created to address a key issue in image aesthetics assessment (IAA): different image themes follow distinct evaluation criteria, yet existing datasets and models generally ignore this fact. TAD66K contains approximately 66,327 carefully curated images covering 47 popular themes, organized into seven super-themes (plants, animals, artifacts, colors, humans, landscapes, and others). Each image was selected manually to ensure thematic clarity and diversity. Every image is annotated with an aesthetic score from 1 to 10, and unlike previous datasets that measure aesthetics, TAD66K provides theme-specific evaluation criteria. For instance, the aesthetic standards applied to images of trees differ from those for rivers or people. 












































\begin{figure*}[t]
  \centering
  % Replace the filename below with your local path if needed
  \includegraphics[width=\linewidth]{magicbrush.jpg}
  \caption{Illustration of \textsc{MagicBrush}: instruction–guided real image editing with both \emph{multi–turn} (left) and \emph{single–turn} (right) workflows. Each triplet consists of a source image, a natural–language instruction, and a target image.}
  \label{fig:magicbrush}
\end{figure*}


\textbf{MagicBrush} is a manually annotated dataset for instruction–guided real image editing. It contains 5,313 edit sessions and 10,388 turns (~10K triplets: source image from MS COCO, instruction, target image). Images are modified with DALL·E~2 until an instruction-consistent, photorealistic target was obtained. As shown in Figure~\ref{fig:magicbrush} some sessions are single-turn, others are multi-turn. 










































\begin{figure}[H]
    \centering
    \includegraphics[width=0.35\textwidth]{pickapic2.jpg}
    \caption{\textbf{Pick-a-Pic} web-app outputs. For each prompt the non-preferred image (left) is darkened and the user-preferred image (right) is highlighted.}
\end{figure}

\textbf{Pick-a-Pic} is a dataset of human preferences for text-to-image generation, collected through a public web application in which users write prompts, generate images with modern diffusion models, and select which image they prefer (or declare a tie). Each datapoint contains a prompt, two model outputs, and a preference label. The initial release contains over $\sim$ 500K pairwise judgments drawn from $\sim$ 35K prompts (and continues to grow), with images produced by Stable Diffusion~2.1, Dreamlike Photoreal~2.0, and SDXL variants under multiple classifier-free guidance (CFG) scales. 


















\begin{figure*}[t]
    \centering
    \includegraphics[width=\textwidth]{internvid.jpg}
    \caption{Examples from \textbf{InternVid}. For each video clip (three frames shown), the dataset provides LLM-generated captions that emphasize nouns (blue) and verbs (green). Compared with raw ASR transcripts (grey), the captions are more visually grounded and descriptive.}
\end{figure*}

\textbf{InternVid} comprises \textasciitilde7.1M public YouTube videos spanning \textasciitilde760K hours and yields \textasciitilde234M trimmed clips (median \textasciitilde10\,s) paired with \textasciitilde4.1B caption words. Coverage is intentionally broad: 16 topical domains, \textasciitilde6,100 action phrases used as queries, and geographically diverse sources; \textasciitilde85\% of videos are 720p (remainder 360–512p). 
































\begin{figure*}[t]
    \centering
    \includegraphics[width=\textwidth]{openvid-1m.jpg}
    \caption{Comparison between existing text-to-video datasets and OpenVid-1M. Datasets such as UCF-101, WebVid-10M, and Panda-70M often contain low-resolution or low-quality videos with short captions, whereas OpenVid-1M features over one million high-quality clips with expressive, detailed textual descriptions highlighting both semantic and visual nuances.}
    \label{fig:openvid1m}
\end{figure*}

\textbf{OpenVid-1M} is a large-scale open-scenario dataset containing over one million high-quality text–video pairs, each accompanied by very descriptive captions, as shown in Figure~\ref{fig:openvid1m}. Prior datasets such as WebVid-10M and Panda-70M prioritize size over fidelity and contain numerous low-resolution, watermarked, flickering, or blurry videos, such is the issue OpenVid-1M seeks to solve. 

The dataset was built by processing large collections such as Panda-50M, ChronoMagic, CelebV-HQ, and Open-Sora-Plan. Videos were evaluated according to multiple criteria, such as using the LAION Aesthetics Predictor to retain only the top 20\% of visually appealing clips. The rather descriptive captions were then generated or rewritten via the multimodal model LLaVA-v1.6-34B.

To further advance high-definition video generation, 433K of these clips at 1080p resolution were selected to create OpenVidHD-0.4M. Overall, OpenVid-1M comprises 1,019,957 clips averaging 7.2 seconds each (approximately 2,051 total hours of content), all at least $512\times512$ in resolution. 

Experimental results confirm that OpenVid-1M and MVDiT significantly outperform prior approaches. Models trained on OpenVid-1M achieve higher aesthetic (VQAA) and technical (VQAT) scores, improved text–video alignment (SD\_score and Blip\_BLEU), and better temporal coherence (Clip\_temp\_score). \\































\begin{figure*}[t]
  \centering
  \includegraphics[width=\textwidth]{visual cot.jpg}
  \caption{\textbf{Visual CoT} examples across five domains. Each question–answer pair is accompanied by an intermediate bounding box (red) that highlights the key region used in reasoning.}
\end{figure*}

\textbf{Visual CoT} contains 438K question–answer pairs, each annotated with a grounding bounding box that marks the critical image region required to answer the question; approximately 98K pairs additionally include step-by-step visual–textual reasoning. Data span five domains: text/doc understanding, fine-grained recognition, charts, general VQA, and relation reasoning, and are built from twelve source corpora (e.g., TextVQA, DocVQA, TextCaps, Flickr30k,GQA, Open Images). \\





























\begin{figure*}[t]
    \centering
    \includegraphics[width=\textwidth]{hq-edit.jpg}
    \caption{(a)–(d): example image edit pairs and instructions from HQ-Edit. (e): comparison of dataset quality between HQ-Edit and previous benchmarks. “Alignment” and “Coherence” are newly proposed metrics introduced in Section~3.4 to evaluate semantic and aesthetic quality.}
    \label{fig:hqedit}
\end{figure*}

\textbf{HQ-Edit} comprises approximately 200,000 instruction-based image edits. Unlike earlier datasets that rely on people to give feedback, HQ-Edit collects data automatically using advanced models: GPT-4V and DALL-E 3. This method creates clearer, higher-quality images and more accurate text instructions, making the link between the text and image much stronger.

In addition to automatic evaluation, a human study on 1,651 image pairs generated by DALL-E~3 was conducted to validate the proposed metrics. Participants rated how well the output images reflected the corresponding edit instructions on a five-point scale, ranging from “totally unrelated” to “perfectly aligned.” Pearson correlation analysis revealed that HQ-Edit’s Alignment metric strongly correlates with human judgments and significantly outperforms the commonly used CLIP Directional Similarity score, which often fails to capture nuanced, instruction-level edits. Furthermore, fine-tuning InstructPix2Pix on HQ-Edit results in a substantial performance boost, achieving a $12.3$-point increase in Alignment and a $5.64$-point improvement in Coherence over the vanilla model.\\













\textbf{BLIP3o-60k} is a 60,000-sample instruction-tuning dataset built by prompting GPT-4o with textual inputs covering various visual domains, including scenes, everyday objects, artistic styles, lighting conditions, gestures, and compositions. Each example in BLIP3o-60k consists of an instruction prompt, a corresponding image, and the image’s CLIP feature embedding. The prompts are designed to mimic real-world creative instructions, such as “a person jumping in the rain wearing a yellow coat” or “an ancient city at sunset in watercolor style”.\\





























\begin{figure*}[t]
  \centering
  \includegraphics[width=\textwidth]{imgedit.jpg}
  \caption{\textbf{ImgEdit} illustrates single–turn edits (add, remove, replace, alter, background/style change, motion change, object extraction, visual and hybrid edits) and multi–turn behaviors (content memory, content understanding, version backtracking).}
\end{figure*}

The \textbf{ImgEdit} dataset contains 1.2M carefully curated edit pairs at high resolution (shorter side $\ge\!1280$), including 1.1M single–turn samples spanning 10 representative operations (local: add, remove, replace, alter, motion change, object extraction; global: background/style; visual edit via reference; hybrid edits that compose multiple local ops) and 110K multi–turn samples covering content memory, content understanding, and version backtracking. 

The original images come from a high–aesthetics subset of LAION (aesthetic score $>\!4.75$ and resolution filtering). These go through multiple steps to generate the edits, including: the creation of concise captions and candidate edit targets with \textbf{GPT-4o}, validation of regions to be edited by CLIPScore and area thresholds (e.g., background edits must occupy $>\!40\%$ of the image), and prompt-guided post-filtering, returning quality scores.\\


























\begin{figure*}[t]
  \centering
  \includegraphics[width=\textwidth]{sharegpt4o.jpg}
  \caption{\textbf{ShareGPT-4o-Image}. Examples from the two splits: \emph{45K} text-to-image (left) and \emph{46K} text\&image-to-image (right) synthesized by GPT-4o-Image.}
\end{figure*}

\textbf{ShareGPT-4o-Image} contains 91K synthetic samples comprising 45K text-to-image pairs and 46K instruction-guided text\&image-to-image triplets, where every target image is produced by GPT-4o-Image. The two distinct sets go through specific pipelines to arrive at the result.

\end{multicols}

\end{document}